\begin{afterword}
\summaryhead

The post quotes Robert Wright's summary of a 1989 Canadian study by Crawford, Salter and Jang. Adults judged how much grief a parent would feel at the death of children of various ages. The resulting grief curve correlated .64 with the reproductive value of children at each age computed from Canadian data, and .92 with the same curve for the !Kung, a hunter-gatherer people. The author, who read the paper after posting, reports N=221.

The author grants that the study asked adults to imagine a parent's grief rather than asking parents, and says other studies agree. The experiment's merits, in the author's view, are these: strong selection could explain so high a correlation; it answers the claim that evolutionary psychology makes no advance predictions; and it shows grief to be an adaptation that runs without tracking present conditions. Grief is ``not even subconsciously'' about reproductive value, or it would follow the Canadian curve. Parents care about children for their own sake. Imagined grief is the right measure, since it is what steers parents before a loss, and the curve that fits is future reproductive value, not past investment. The post ends: humans and natural selection are insane in different ways.

\responsehead

The study is reported accurately as far as it can be checked. Wright's summary matches the book it comes from, \textsc{The Moral Animal}, including the .64 and the .92, and the author read the paper and reported that the correlations were as described. The post also raises, unprompted, the question whether today's !Kung are typical of ancestral hunter-gatherers.

The problem is the weight put on .92. By Wright's description, it is a correlation between two curves, one point per age of child; a commenter who read the paper counted ten ages. The 221 in ``correlation .92 and N=221 is pretty strong evidence'' are raters, whose answers were averaged into those points. On ten points .92 is still unlikely by chance, but it is ten points. A correlation of averages over hundreds of raters can come out near 1 even when individuals disagree widely, so it is not comparable with correlations between individuals' answers. The post's comparison, ``one of the highest correlations I've seen in any psychology experiment,'' sets it beside all of them.

The conclusions go further than the comparison. That grief is ``not even subconsciously'' about reproductive value rests on the gap between .92 and .64, computed on the same few points; the post does not say whether that gap could be chance. That parents ``care about children for their own sake'' is not something the study measured. Robin Hanson made this objection in the comments, and the author replied, ``I wasn't arguing for the other pole.''

The defence of imagined grief fits the study only indirectly. The post argues that imagined grief is what steers parents before a loss. The subjects were not imagining their own loss; they estimated what a parent would feel. For the link to real parents the post relies on its ``understanding'' of other studies, without naming one. Two later studies of bereaved parents, cited in a 2020 review, found grief highest for children lost in adolescence, which supports the post here.

In short: an accurately reported study, whose .92 is a correlation between averages at about ten ages rather than evidence from 221 observations, and whose conclusion about what parents care about goes beyond what it tested.
\end{afterword}
